% !TEX root = ../main.tex
% ============================================================================
%                        HEADER & FOOTER STYLE
% ============================================================================
% Nomor halaman: kanan atas (untuk halaman biasa)
%                tengah bawah (untuk awal bab/chapter)
\pagestyle{fancy}
\fancyhf{}
\fancyhead[R]{\thepage}  % Nomor halaman di kanan atas
\renewcommand{\headrulewidth}{0pt}
\setlength{\headheight}{14.5pt}

% Gaya halaman awal bab (plain) -> nomor di tengah bawah
\fancypagestyle{plain}{%
  \fancyhf{}
  \fancyfoot[C]{\thepage}
  \renewcommand{\headrulewidth}{0pt}
  \renewcommand{\footrulewidth}{0pt}
}

% ============================================================================
%                        FORMAT JUDUL BAB & SUBBAB
% ============================================================================
% Penomoran Chapter dengan Romawi (BAB I, BAB II, dst)
\renewcommand{\thechapter}{\Roman{chapter}}

% Font untuk heading utama (14pt, bold)
\newcommand{\HeadingFont}{\fontsize{14pt}{21pt}\selectfont\bfseries}

% ==== FORMAT CHAPTER (BAB) ====
% Contoh hasil: "BAB I" (baris 1), "PENDAHULUAN" (baris 2)
\titleformat{\chapter}[display]
  {\normalfont\HeadingFont\centering}  % Format: 14pt, bold, tengah
  {\MakeUppercase{BAB \thechapter}}    % Label: BAB I, BAB II, dst
  {0pt}                                % Jarak antara label dan judul
  {\HeadingFont\MakeUppercase}         % Format judul: uppercase, bold, 14pt

% Spacing Chapter: {left}{before-sep}{after-sep}
\titlespacing*{\chapter}
  {0pt}              % left margin
  {0pt}              % space before (dari atas halaman)
  {1.5\baselineskip} % space after (ke konten berikutnya)

% ==== FORMAT SECTION (Heading 1: 1.1, 1.2, dst) ====
\renewcommand{\thesection}{\arabic{chapter}.\arabic{section}}
\titleformat{\section}[block]
  {\normalfont\bfseries}  % Format: bold, 12pt
  {\thesection}           % Nomor section
  {1em}                   % Jarak antara nomor dan judul
  {}                      % Kode sebelum judul

% Spacing Section
\titlespacing*{\section}
  {0pt}              % left margin
  {1\baselineskip}   % space before
  {0.5\baselineskip} % space after

% ==== FORMAT SUBSECTION (Heading 2: 1.1.1, 1.1.2, dst) ====
\renewcommand{\thesubsection}{\arabic{chapter}.\arabic{section}.\arabic{subsection}}
\titleformat{\subsection}[block]
  {\normalfont\bfseries}  % Format: bold, 12pt
  {\thesubsection}        % Nomor subsection
  {1em}                   % Jarak antara nomor dan judul
  {}                      % Kode sebelum judul

% Spacing Subsection
\titlespacing*{\subsection}
  {0pt}              % left margin
  {1\baselineskip}   % space before
  {0.5\baselineskip} % space after

% ==== FORMAT SUBSUBSECTION (Heading 3: 1.1.1.1, dst) ====
\renewcommand{\thesubsubsection}{\arabic{chapter}.\arabic{section}.\arabic{subsection}.\arabic{subsubsection}}
\titleformat{\subsubsection}[block]
  {\normalfont\bfseries}  % Format: bold, 12pt
  {\thesubsubsection}     % Nomor subsubsection
  {1em}                   % Jarak antara nomor dan judul
  {}                      % Kode sebelum judul

% Spacing Subsubsection
\titlespacing*{\subsubsection}
  {0pt}              % left margin
  {1\baselineskip}   % space before
  {0.5\baselineskip} % space after

% ==== FORMAT PARAGRAPH (Heading 4) ====
\setcounter{secnumdepth}{4}  % Aktifkan penomoran hingga paragraph
\titleformat{\paragraph}[runin]
  {\normalfont\bfseries}  % Format: bold, inline dengan teks
  {\theparagraph}         % Nomor paragraph
  {1em}                   % Jarak antara nomor dan judul
  {}                      % Kode sebelum judul
  [.]                     % Tanda setelah judul (titik)

% Spacing Paragraph
\titlespacing*{\paragraph}
  {0pt}              % left margin
  {1\baselineskip}   % space before
  {0.5em}            % space after (horizontal karena runin)

% ============================================================================
%                          INDENTASI PARAGRAF
% ============================================================================
% Indent ~4 karakter (1.5em untuk Times 12pt)
\setlength{\parindent}{1.5em}

% ============================================================================
%                   JUSTIFY & HYPHENATION
% ============================================================================
\setlength{\emergencystretch}{1.5em}
\AtBeginEnvironment{verbatim}{\singlespacing}

% ============================================================================
%                 COMMAND RENDER HALAMAN DAFTAR
% ============================================================================
\makeatletter
\newcommand{\RenderListPage}[3][]{%
  \cleardoublepage
  \phantomsection
  \thispagestyle{plain}%
  \setstretch{1.5}%
  #1%
  {\HeadingFont\centering \MakeUppercase{#2}\par}\vspace{\baselineskip}%
  \noindent\hfill\textbf{Halaman}\par\vspace{0.5\baselineskip}%
  \begingroup\setstretch{1.5}\@starttoc{#3}\endgroup\par}

\newcommand{\RenderListPageNoSpace}[3][]{%
  \cleardoublepage
  \phantomsection
  \thispagestyle{plain}%
  \setstretch{1.5}%
  #1%
  {\HeadingFont\centering \MakeUppercase{#2}\par}\vspace{\baselineskip}%
  \noindent\hfill\textbf{Halaman}\par\vspace{0.5\baselineskip}%
  \begingroup
  \setstretch{1.5}
  \let\addvspace\@gobble
  \@starttoc{#3}%
  \endgroup\par}
\makeatother

% ============================================================================
%                         KONFIGURASI DAFTAR ISI (TOC)
% ============================================================================
\setcounter{tocdepth}{3}
\setcounter{secnumdepth}{3}

\renewcommand{\cftchappresnum}{BAB~}
\renewcommand{\cftchapaftersnum}{}
\setlength{\cftchapindent}{0em}
\setlength{\cftchapnumwidth}{4em}

\renewcommand{\cftchapfont}{\bfseries}
\renewcommand{\cftchappagefont}{\bfseries}
\renewcommand{\cftchapleader}{\hfill} % Menghapus titik-titik

\renewcommand{\cftsecfont}{\normalfont}
\renewcommand{\cftsecpagefont}{\normalfont}
\renewcommand{\cftsecleader}{\hfill} % Menghapus titik-titik
\setlength{\cftsecindent}{2em}
\setlength{\cftsecnumwidth}{2em}

\renewcommand{\cftsubsecleader}{\hfill} % Menghapus titik-titik
\setlength{\cftsubsecindent}{4em}
\setlength{\cftsubsecnumwidth}{3em}

\renewcommand{\cftsubsubsecleader}{\hfill} % Menghapus titik-titik
\setlength{\cftsubsubsecindent}{7em}
\setlength{\cftsubsubsecnumwidth}{4em}

\setlength{\cftbeforechapskip}{0pt}
\setlength{\cftbeforesecskip}{0pt}
\setlength{\cftbeforesubsecskip}{0pt}
\setlength{\cftbeforesubsubsecskip}{0pt}

\renewcommand{\cftdotsep}{1}

\addto\captionsindonesian{%
  \renewcommand{\contentsname}{Daftar Isi}
  \renewcommand{\listfigurename}{Daftar Gambar}
  \renewcommand{\listtablename}{Daftar Tabel}
  \renewcommand{\bibname}{Daftar Pustaka}
  \renewcommand{\refname}{Daftar Pustaka}
  \renewcommand{\appendixname}{LAMPIRAN}
}

\renewcommand{\cfttoctitlefont}{\hfill\HeadingFont}
\renewcommand{\cftaftertoctitle}{\hfill}
\renewcommand{\cftloftitlefont}{\hfill\HeadingFont}
\renewcommand{\cftafterloftitle}{\hfill}
\renewcommand{\cftlottitlefont}{\hfill\HeadingFont}
\renewcommand{\cftafterlottitle}{\hfill}

\setlength{\cftbeforetoctitleskip}{0pt}
\setlength{\cftaftertoctitleskip}{1\baselineskip}
\setlength{\cftbeforeloftitleskip}{0pt}
\setlength{\cftafterloftitleskip}{1\baselineskip}
\setlength{\cftbeforelottitleskip}{0pt}
\setlength{\cftafterlottitleskip}{1\baselineskip}

% ==== Global: Hapus Titik-titik pada SEMUA daftar (LoF, LoT, dll) ====
\makeatletter
\renewcommand{\@dottedtocline}[5]{%
  \ifnum #1>\c@tocdepth \else
    \vskip \z@ \@plus.2\p@
    {\leftskip #2\relax \rightskip \@tocrmarg \parfillskip -\rightskip
     \parindent #2\relax\@afterindenttrue
     \interlinepenalty\@M
     \leavevmode
     \@tempdima #3\relax
     \advance\leftskip \@tempdima \null\nobreak\hskip -\leftskip
     {#4}\nobreak
     \hfill % Mengganti \leaders\hbox{$\m@th\mkern \@dotsep mu.\mkern \@dotsep mu$}\hfill
     \nobreak
     \hb@xt@\@pnumwidth{\hfil\normalfont \normalcolor #5}%
     \par}%
  \fi}
\makeatother

% ==== Standarisasi Daftar Gambar (LoF) ====
\setlength{\cftfignumwidth}{3.5em}
\setlength{\cftfigindent}{0em}
\setlength{\cftbeforefigskip}{0pt}
\renewcommand{\cftfigleader}{\hfill} % Menghapus titik-titik
\renewcommand{\cftfigpresnum}{}
\renewcommand{\cftfigaftersnum}{\space}

% ==== Standarisasi Daftar Tabel (LoT) ====
\setlength{\cfttabnumwidth}{3.5em}
\setlength{\cfttabindent}{0em}
\setlength{\cftbeforetabskip}{0pt}
\renewcommand{\cfttableader}{\hfill} % Menghapus titik-titik
\renewcommand{\cfttabpresnum}{}
\renewcommand{\cfttabaftersnum}{\space}

% ============================================================================
%                           DAFTAR PERSAMAAN
% ============================================================================
\newcommand{\listequationname}{Daftar Persamaan}
\newcommand{\eqnumberwithparen}{(\arabic{chapter}.\arabic{equation})}
\newlistof{myequations}{equ}{\listequationname}
\newcommand{\myequations}[1]{%
  \addcontentsline{equ}{myequations}{\protect\numberline{\eqnumberwithparen\enspace}#1}\ignorespaces}
\setlength{\cftmyequationsnumwidth}{4em}
\renewcommand{\cftmyequationsleader}{\hfill} % Menghapus titik-titik
\providecommand{\cftmyequationstitlefont}{}
\providecommand{\cftaftermyequationstitle}{}
\renewcommand{\cftmyequationstitlefont}{\hfill\HeadingFont}
\renewcommand{\cftaftermyequationstitle}{\hfill}
\newcommand{\printlistequations}{%
  \RenderListPage[\addcontentsline{toc}{chapter}{\MakeUppercase{\listequationname}}]{\listequationname}{equ}}

% ============================================================================
%                         DAFTAR ALGORITMA/KODE
% ============================================================================
\newcommand{\listalgorithmname}{Daftar Kode dan Algoritma}
\providecommand{\cftloaleader}{\hfill} % Menggunakan providecommand agar aman
\renewcommand{\cftloaleader}{\hfill}
\newcommand{\printlistalgorithms}{%
  \RenderListPageNoSpace[\addcontentsline{toc}{chapter}{\MakeUppercase{\listalgorithmname}}]{\listalgorithmname}{loa}%
}

% ============================================================================
%                           DAFTAR LAMPIRAN
% ============================================================================
\newcommand{\listappendixname}{Daftar Lampiran}
\newlistof{myappendices}{apx}{\listappendixname}
\newcommand{\myappendix}[2]{%
\addcontentsline{apx}{myappendices}{\protect\numberline{Lampiran #1\enspace}#2}\par}
\setlength{\cftmyappendicesnumwidth}{7em}
\renewcommand{\cftmyappendicesleader}{\hfill} % Menghapus titik-titik
\providecommand{\cftmyappendicestitlefont}{}
\providecommand{\cftaftermyappendicestitle}{}
\renewcommand{\cftmyappendicestitlefont}{\hfill\HeadingFont}
\renewcommand{\cftaftermyappendicestitle}{\hfill}
\newcommand{\printlistappendices}{%
  \RenderListPage[\addcontentsline{toc}{chapter}{\MakeUppercase{\listappendixname}}]{\listappendixname}{apx}}

% ============================================================================
%                      DAFTAR ISTILAH (GLOSARIUM)
% ============================================================================
\newcommand{\listtermname}{Daftar Istilah}
\newcommand{\myterm}[2]{#1 & -- & #2 \\} % Default placeholder
\newcommand{\printlistterms}{%
  \cleardoublepage
  \phantomsection
  \thispagestyle{plain}%
  \setstretch{1.5}%
  \addcontentsline{toc}{chapter}{\MakeUppercase{\listtermname}}%
  {\HeadingFont\centering \MakeUppercase{\listtermname}\par}\vspace{\baselineskip}%
  \noindent
  \begingroup
  \setstretch{1.5}
  \renewcommand{\myterm}[2]{##1 & -- & ##2 \\[0.5em]}
  \begin{xltabular}{\textwidth}{@{} >{\raggedright\arraybackslash}p{3.5cm} @{} c @{\hspace{1em}} X @{}}
    \input{frontmatter/daftar_istilah}
  \end{xltabular}
  \endgroup
}

% ============================================================================
%                          GAYA CAPTION INDONESIA
% ============================================================================
\captionsetup[figure]{
  name=Gambar,
  labelsep=colon,
  font=small,
  position=bottom,
  skip=6pt
}

\captionsetup[table]{
  name=Tabel,
  labelsep=colon,
  font=small,
  position=top,
  skip=6pt,
  justification=centering
}

% ============================================================================
%                      ALGORITMA: PENOMORAN & NAMA
% ============================================================================
\renewcommand{\thealgocf}{\arabic{chapter}.\arabic{algocf}}
\makeatletter
\@addtoreset{algocf}{chapter}
\makeatother

\SetAlgorithmName{Algoritma}{Algoritma}{DAFTAR KODE DAN ALGORITMA}

% Spasi 1.0 di dalam body algoritma/pseudocode dan listing
\AtBeginEnvironment{algorithm}{\singlespacing}
\AtBeginEnvironment{lstlisting}{\singlespacing}

% ============================================================================
%                         KODE SUMBER (LISTINGS)
% ============================================================================
\makeatletter
\AtBeginDocument{
  \@addtoreset{algocf}{chapter}
  \@addtoreset{lstlisting}{chapter}
  \renewcommand{\thelstlisting}{\thealgocf}
  \let\l@lstlisting\l@algocf
}
\makeatother

\renewcommand{\lstlistingname}{Kode}
\renewcommand{\lstlistlistingname}{DAFTAR KODE SUMBER}

% ============================================================================
%                         TABEL: KONFIGURASI
% ============================================================================
\sisetup{
  detect-weight=true,
  group-separator={,},
  group-minimum-digits=4,
  table-format=3.1,
}

\setlength{\arrayrulewidth}{0.3pt}
\renewcommand{\arraystretch}{1.15}

% ============================================================================
%                     SISTEM PENOMORAN (IEEE STYLE)
% ============================================================================
\numberwithin{figure}{chapter}
\renewcommand{\thefigure}{\arabic{chapter}.\arabic{figure}}

\numberwithin{table}{chapter}
\renewcommand{\thetable}{\arabic{chapter}.\arabic{table}}

\numberwithin{equation}{chapter}
\renewcommand{\theequation}{\arabic{chapter}.\arabic{equation}}

% ============================================================================
%                        FORMAT LAMPIRAN
% ============================================================================
\renewcommand{\appendixname}{LAMPIRAN}

\makeatletter
\g@addto@macro\appendix{%
  \renewcommand{\chaptername}{\appendixname}%
  \renewcommand{\thechapter}{\Alph{chapter}}%
  \titleformat{\chapter}[display]
  {\normalfont\fontsize{14pt}{16pt}\selectfont\bfseries\centering}
  {\MakeUppercase{\appendixname~\thechapter}}
  {0pt}
  {\MakeUppercase}
}
\makeatother
